\chapter{Method}
\label{ch:method}

\section{Dataset and Output Format}

The model predicts a two-dimensional output vector:

\begin{equation}
  y = [y_x, y_y].
\end{equation}

The evaluation compares predicted outputs against the ground-truth labels \texttt{y\_true.npy}.

\section{Python Export}

Python export produces the deployment assets:

\begin{itemize}
  \item \texttt{export\_ir/export.json}: model structure and file references.
  \item \texttt{*.npy}: weights, biases, and quantization scales.
  \item \texttt{float/samples.npy}: input samples.
  \item \texttt{float/y\_true.npy}: ground-truth labels.
  \item \texttt{float/y\_float.npy}: PyTorch float model outputs.
\end{itemize}

\section{C++ Reference Inference}

The C++ reference reads the exported assets and runs fake, int8, or int16 deployment-oriented inference. Its output is stored as:

\begin{lstlisting}
final_hw/cases/<case>/float/cpp_full_int16.npy
\end{lstlisting}

This output defines the deployment quantization reference. It is not a cycle-accurate or stage-accurate RTL golden model.

\section{Python Hardware-Like Simulation}

The Python \texttt{hw\_like} path reproduces RTL-oriented fixed-point behavior using NumPy. It models Q8.8 activations, Q0.16 sigmoid outputs, Q1.15 lambda values, rounding, saturation, memory layout, and stage boundaries.

\section{Python C++-Like Trace}

The Python \texttt{cpp\_like} path reproduces the C++ reference semantics in Python. It is used to obtain intermediate stage traces for debug. It does not call the C++ executable.

\section{RTL Asset Generation}

The hardware-like path generates:

\begin{itemize}
  \item Weight memory files.
  \item Scale memory files.
  \item Stage golden files such as \texttt{ssm\_golden\_q88.mem} and \texttt{gate\_y\_golden\_q88.mem}.
  \item Control files for the RTL testbench.
\end{itemize}

\section{Error Metrics}

For final two-dimensional predictions, the Euclidean error is:

\begin{equation}
  e_i = \sqrt{(y_{i,x} - y^{\ast}_{i,x})^2 + (y_{i,y} - y^{\ast}_{i,y})^2}.
\end{equation}

The reported metrics include mean error, median error, percentile error, and maximum error.

For tensor comparisons, mean absolute error is:

\begin{equation}
  \mathrm{MAE}(a,b) = \frac{1}{N}\sum_{i=1}^{N}|a_i - b_i|.
\end{equation}

The maximum absolute error is:

\begin{equation}
  \mathrm{max\_abs}(a,b) = \max_i |a_i - b_i|.
\end{equation}

\section{Metric Interpretation}

\begin{table}[h]
\centering
\caption{Interpretation of key comparisons.}
\label{tab:metric-meaning}
\begin{tabular}{lll}
\toprule
Comparison & Purpose & Large value indicates \\
\midrule
\texttt{float vs y\_true} & Original model quality & Model accuracy limitation \\
\texttt{cpp vs float} & Quantization loss & Poor quantization strategy \\
\texttt{hw\_like vs cpp} & Hardware extra error & Q format, scale, rounding, or stage mismatch \\
\texttt{stage MAE} & Stage localization & First stage where hardware semantics diverge \\
\bottomrule
\end{tabular}
\end{table}

